\section{Holonomy classification}

The path-receipt construction respects concatenation and cancellation of a
dart with its reverse. It therefore descends from based loops to the graph
fundamental group.

\begin{theorem}[Gauge classification]
Let $X$ be a connected finite graph and $R$ a finite group. Gauge-equivalence
classes of inverse-compatible regular $R$-receipt lifts of $X$ are in
bijection with $R$-conjugacy classes of homomorphisms
\[
\rho:\pi_1(X,r)\longrightarrow R.
\]
\end{theorem}

\begin{proof}
A gauge transformation conjugates every based loop receipt by the single
element $g(r)$, so gauge-equivalent assignments give conjugate
representations.

Conversely, fix a spanning tree and normalize two assignments on it. They are
then determined by their $\beta$ cotree receipts. If their representations
are conjugate by $t\in R$, the constant gauge $g(v)=t$ simultaneously
conjugates every cotree receipt while preserving identity on the tree.
Composing this constant gauge with the two normalizing gauges proves gauge
equivalence.
\end{proof}

Changing the root transports the representation by a path receipt and hence
changes it only by conjugacy. Changing the spanning tree changes the free
generating basis of $\pi_1(X,r)$ but not the underlying representation.

When $R$ is abelian, conjugacy is trivial and a gauge class is an element of
\[
\operatorname{Hom}(\pi_1(X,r),R)\cong R^{\beta}.
\]
For nonabelian $R$, simultaneous conjugacy is essential: individual circuit
values are coordinate-dependent, while the representation class is not.
